% !Mode:: "TeX:UTF-8"
% ============================================================
% PREAMBLE: all packages and settings
% Compiler: XeLaTeX (latexmk selects it automatically)
% ============================================================

% ------------------------------------------------------------
% 1. MATHEMATICS (loaded BEFORE unicode-math)
% ------------------------------------------------------------
\usepackage{amsmath}
\usepackage{amsthm}
\usepackage{mathtools}
% tikz is not loaded (it is slow): figures are ready-made PDFs, their sources are in rasmlar/manba/
% Bold vector: \vect{v}.

% ------------------------------------------------------------
% 2. FONTS
% ------------------------------------------------------------
% Fonts are loaded by FILE NAME: searching the system fonts
% (fontconfig) is very slow on Overleaf, while a file name is found at once.
% NOTE (for speed): the font's "TeX ligatures" are switched off. Type the
% final characters in the text:  ’ (Cramer’s),  – (en dash),  — ,  “ ” (quotes).
% (If `--` or ``...'' is typed, it is printed unchanged.)
\usepackage{amssymb}
% Both XeLaTeX and PDFLaTeX work (for Overleaf and WinEdt):
\usepackage{iftex}
\ifXeTeX
  % newtxmath, Times-style formulas (several times faster than unicode-math,
  % and matching the text font). Loaded BEFORE fontspec.
  \usepackage{newtxmath}
  \usepackage[no-math]{fontspec}
  % Main font: TeX Gyre Termes (Overleaf, TeX Live, MiKTeX package "tex-gyre").
  % If it is not found (e.g. the package is not installed in MiKTeX), Windows'
  % Times New Roman is used, and failing that, Latin Modern.
  \IfFontExistsTF{texgyretermes-regular.otf}
    {\setmainfont{texgyretermes}[Extension=.otf, Ligatures=TeXOff,
        UprightFont=*-regular, BoldFont=*-bold,
        ItalicFont=*-italic, BoldItalicFont=*-bolditalic]}
    {\IfFontExistsTF{Times New Roman}
      {\setmainfont{Times New Roman}[Ligatures=TeXOff]}
      {\setmainfont{Latin Modern Roman}[Ligatures=TeXOff]}}
  % Digits and names such as sin, cos in formulas are also set in Times:
  \DeclareSymbolFont{operators}{OT1}{ntxtlf}{m}{n}
  \SetSymbolFont{operators}{bold}{OT1}{ntxtlf}{b}{n}
  \DeclareMathAlphabet{\mathrm}{OT1}{ntxtlf}{m}{n}
  \DeclareMathAlphabet{\mathbf}{OT1}{ntxtlf}{b}{n}
  \DeclareMathAlphabet{\mathit}{OT1}{ntxtlf}{m}{it}
  \providecommand{\symbf}[1]{\mathbf{#1}}
\else
  % ---- PDFLaTeX (the "PDFLaTeX" button in WinEdt) ----
  \usepackage[T2A,T1]{fontenc}       % T1 -- Latin, T2A -- Cyrillic letters
  \usepackage[utf8]{inputenc}
  \usepackage{newtxtext}             % Times-style text font
  \usepackage{newtxmath}             % Times-style formulas
  \providecommand{\symbf}[1]{\mathbf{#1}}
\fi

% ------------------------------------------------------------
% 3. LANGUAGE
% ------------------------------------------------------------
\usepackage[english]{babel}

% ------------------------------------------------------------
% 4. PAGE
% ------------------------------------------------------------
\usepackage[
    a5paper,
    inner=2.0cm,
    outer=1.6cm,
    top=1.8cm,
    bottom=1.8cm,
    bindingoffset=0.3cm,
    headheight=16pt,
    headsep=0.5cm,
    footskip=0.8cm
]{geometry}

\usepackage{setspace}
\setstretch{1.1}
\setlength{\parindent}{1em}
\setlength{\parskip}{0pt}
% microtype is not used because it slows down compilation
\setlength{\emergencystretch}{3em}
\tolerance=1000
% Automatic line breaking of a wide display made of several formulas:
%   \[\kf{ f_1,\qquad f_2 }\]  -- on one line if it fits; otherwise
% it breaks at a \qquad (\quad) or after an "=" or "\to" sign.
\newcommand{\kfqquad}{\penalty0\hskip2em plus 1em\relax}
\newcommand{\kfquad}{\penalty100\hskip1em plus .5em\relax}
\newcommand{\kf}[1]{%
  \vcenter{\hsize=\linewidth\leftskip=0pt plus 1fil\rightskip=0pt plus 1fil
    \parfillskip=0pt\parindent=0pt\emergencystretch=0pt
    \relpenalty=300 \binoppenalty=900
    \noindent\strut$\displaystyle\let\qquad\kfqquad\let\quad\kfquad #1$\strut\par}}
\clubpenalty=10000
\widowpenalty=10000

% ------------------------------------------------------------
% 5. COLOURS, FRAMES, THEOREM ENVIRONMENTS
% ------------------------------------------------------------
\usepackage{xcolor}
\usepackage{etoolbox}
\usepackage{framed}
% !Mode:: "TeX:UTF-8"
% ============================================================
% DEFINITION, THEOREM, EXAMPLE AND OTHER ENVIRONMENTS
% Numbered within sections: Definition 1.2.3, Example 1.2.4 ...
% ============================================================

\definecolor{tarifrang}{RGB}{40,70,140}
\definecolor{tariffon}{RGB}{244,247,253}
\definecolor{misolrang}{RGB}{175,110,30}
\definecolor{misolfon}{RGB}{253,248,238}
\definecolor{teoremarang}{RGB}{90,90,90}
\definecolor{teoremafon}{RGB}{246,246,246}

% Frames use the "framed" package: much faster than mdframed, and
% they break automatically at the end of a page.
\newenvironment{rangliramka}[2]{%
  \def\FrameCommand{\fboxsep=5pt\fboxrule=0.8pt\fcolorbox{#1}{#2}}%
  \MakeFramed{\advance\hsize-\width\FrameRestore}}%
  {\endMakeFramed}
\newenvironment{chapramka}{%
  \def\FrameCommand{{\color{teoremarang}\vrule width 2pt}\fboxsep=4pt\colorbox{teoremafon}}%
  \MakeFramed{\advance\hsize-\width\FrameRestore}}%
  {\endMakeFramed}

% --- Theorem style (italic) ---
\theoremstyle{plain}
\newtheorem{teorema}{Theorem}[section]
\newtheorem{lemma}[teorema]{Lemma}
\newtheorem{natija}[teorema]{Corollary}
\newtheorem{xossa}[teorema]{Property}

% --- Definition style (upright) ---
\theoremstyle{definition}
\newtheorem{tarif}{Definition}[section]
\newtheorem{misol}{Example}[section]
\newtheorem{mashq}{Exercise}[section]
\newenvironment{izoh}{\smallskip\noindent\textit{Note. }}{\smallskip}
\theoremstyle{remark}
\newtheorem{eslatma}{Remark}[section]

% --- Frames ---
\BeforeBeginEnvironment{tarif}{\begin{rangliramka}{tarifrang}{tariffon}}
\AfterEndEnvironment{tarif}{\end{rangliramka}}
\BeforeBeginEnvironment{misol}{\begin{rangliramka}{misolrang}{misolfon}}
\AfterEndEnvironment{misol}{\end{rangliramka}}
\BeforeBeginEnvironment{teorema}{\begin{chapramka}}
\AfterEndEnvironment{teorema}{\end{chapramka}}
\BeforeBeginEnvironment{lemma}{\begin{chapramka}}
\AfterEndEnvironment{lemma}{\end{chapramka}}

% --- Solution (in place of a proof) ---
\newenvironment{yechish}{\begin{proof}[Solution]}{\end{proof}}
\renewcommand{\qedsymbol}{\rule{0.5em}{0.5em}}
\addto\captionsenglish{\renewcommand{\proofname}{Proof}}


% ------------------------------------------------------------
% 6. FIGURES, TABLES, LISTS, COLUMNS
% ------------------------------------------------------------
\usepackage{graphicx}
\usepackage{fancyvrb}
% Program code: \begin{Verbatim}[fontsize=\footnotesize,frame=single,label={\small Python}] ... \end{Verbatim}
\makeatletter
\providecommand{\LenToUnit}[1]{\strip@pt\dimexpr#1*\p@/\unitlength\relax}
\makeatother
\graphicspath{{rasmlar/}{../rasmlar/}}
\usepackage{booktabs}
\usepackage{array}
\usepackage{caption}
\captionsetup{font=small, labelsep=period}
\usepackage{multicol}
\usepackage{enumitem}
\setlist[itemize]{leftmargin=1.5em}
\setlist[enumerate]{leftmargin=1.8em}
\usepackage{tasks}
\settasks{label=\arabic*., label-width=1.6em, item-indent=2.2em}

% ------------------------------------------------------------
% 7. HEADINGS
% ------------------------------------------------------------
\usepackage{titlesec}
\titleformat{\chapter}[display]
    {\normalfont\bfseries\centering}
    {\Large\chaptertitlename\ \thechapter}
    {10pt}
    {\LARGE\raggedright\centering}
\titleformat{name=\chapter,numberless}[display]
    {\normalfont\bfseries\centering}
    {}
    {0pt}
    {\LARGE}
\titlespacing*{\chapter}{0pt}{10pt}{25pt}

\titleformat{\section}
    {\normalfont\large\bfseries\raggedright}
    {\thesection}{0.6em}{}
\titleformat{\subsection}
    {\normalfont\normalsize\bfseries\raggedright}
    {\thesubsection}{0.6em}{}
\titleformat{\subsubsection}
    {\normalfont\normalsize\bfseries\itshape\raggedright}
    {\thesubsubsection}{0.6em}{}

\setcounter{secnumdepth}{2}   % numbered down to 1.1.1
\setcounter{tocdepth}{1}      % chapters and sections in the contents

% ------------------------------------------------------------
% 8. RUNNING HEADS
% ------------------------------------------------------------
\usepackage{emptypage}        % no running heads on empty pages
\usepackage{truncate}         % a long heading is truncated to one line
\usepackage{fancyhdr}
\pagestyle{fancy}
\fancyhf{}
\fancyhead[LE]{\small\truncate{\headwidth}{\leftmark}}
\fancyhead[RO]{\small\truncate{\headwidth}{\hfill\rightmark}}
\fancyfoot[LE,RO]{\small\thepage}
\renewcommand{\headrulewidth}{0.4pt}
\renewcommand{\footrulewidth}{0pt}
\renewcommand{\chaptermark}[1]{\markboth{\chaptername\ \thechapter. #1}{}}
\renewcommand{\sectionmark}[1]{\markright{\thesection.\ #1}}
% First page of a chapter: page number only
\fancypagestyle{plain}{%
    \fancyhf{}%
    \fancyfoot[LE,RO]{\small\thepage}%
    \renewcommand{\headrulewidth}{0pt}}

% ------------------------------------------------------------
% 9. NAMES (set explicitly over the ones installed by babel)
% ------------------------------------------------------------
\addto\captionsenglish{%
    \renewcommand{\contentsname}{CONTENTS}%
    \renewcommand{\bibname}{REFERENCES}%
    \renewcommand{\figurename}{Figure}%
    \renewcommand{\tablename}{Table}%
    \renewcommand{\chaptername}{Chapter}}

% ------------------------------------------------------------
% 10. SHORT COMMANDS
% ------------------------------------------------------------
\newcommand{\R}{\mathbb{R}}
\providecommand{\C}{}\renewcommand{\C}{\mathbb{C}}
\newcommand{\N}{\mathbb{N}}
\newcommand{\Z}{\mathbb{Z}}
\newcommand{\vect}[1]{\symbf{#1}}
\DeclareMathOperator{\rg}{rg}
\DeclareMathOperator{\tr}{tr}
\numberwithin{equation}{section}
% Section (topic) heading:
%   \mashgulot{1}{The Concept of a Matrix}
% prints a numbered section; the full title goes to the contents.
\makeatletter
\newcommand{\mashgulot}[2]{%
    \section{#2}%
    \expandafter\xdef\csname mavzu@#1\endcsname{\thesection}}
% Heading of the answers to topic #1 (printed with its section number):
\newcommand{\javobmashgulot}[2]{%
  \subsection*{\ifcsname mavzu@#1\endcsname\csname mavzu@#1\endcsname.\ \fi#2}}
\makeatother
% ------------------------------------------------------------
% 11. COMPILING CHAPTERS SEPARATELY
% ------------------------------------------------------------
\usepackage{subfiles}

% ------------------------------------------------------------
% 11a. SINGLE-PASS COMPILATION (for the 10-second limit on Overleaf)
% No .aux file is written, so xelatex runs ONLY ONCE each time.
% The table of contents is written to \jobname.mnd and read in the next
% compilation; if that file does not exist yet, the ready-made
% qismlar/mundarija.mnd is read. (\label/\ref are not used.)
% ------------------------------------------------------------
\nofiles
\makeatletter
\newwrite\mnd@out
\newif\ifmnd@open
\renewcommand\addtocontents[2]{%
  \def\@tempa{#1}\def\@tempb{toc}%
  \ifx\@tempa\@tempb
    \ifmnd@open
      % (\nofiles disables \protected@write, hence our own copy)
      \begingroup\let\thepage\relax
        \let\label\@gobble\let\index\@gobble\let\glossary\@gobble
        \let\protect\@unexpandable@protect
        \edef\reserved@a{\write\mnd@out{#2}}\reserved@a
      \endgroup
      \if@nobreak\ifvmode\nobreak\fi\fi
    \fi
  \fi}
\renewcommand\@starttoc[1]{%
  \begingroup\makeatletter
  \IfFileExists{\jobname.mnd}{\@input{\jobname.mnd}}{\@input{qismlar/mundarija.mnd}}%
  \endgroup
  \immediate\openout\mnd@out=\jobname.mnd\relax
  \global\mnd@opentrue}
\makeatother

% ------------------------------------------------------------
% 12. HYPERLINKS (loaded last)
% ------------------------------------------------------------
% hyperref slows compilation down about twice and is not needed for a
% printed book. If clickable links in the PDF are wanted, remove the % sign
% at the beginning of the following line:
% \usepackage[hidelinks,bookmarksnumbered]{hyperref}
\providecommand{\texorpdfstring}[2]{#1}
\providecommand{\phantomsection}{}
\providecommand{\url}[1]{\texttt{#1}}
